\paragraph{Implications for Idea 1.}
The evidence supports retaining a grounded supervision stage, while narrowing what the current analysis establishes. Released frame-grid predictions outperform description predictions on matched items, including items missed by a text-only probe. However, a potential description-prompt defect prevents interpreting the joint-score gap as a clean measure of perceptual necessity. Before training, we will verify the target backbone's frame handling and rerun matched, consistently prompted controls. We will preserve the original comparison of the curriculum with an equal-composition unordered mixture and order ablations; the present analysis does not test that central hypothesis.

Three data decisions follow. First, deduplicate NormBank and explicitly resolve conflicting context labels before splitting or sampling. Second, preserve NormLens's native judgment task and annotator disagreement rather than generating stylistic multiple-choice distractors. Third, use probe-error results alongside full-set results, with text-only controls and equal temporal evidence for video versus frames. A supervised EgoNormia diagnostic probe is not the transfer model, but it does use benchmark labels: we should describe the benchmark as excluded from transfer-model training, not as untouched by analysis. The team has contacted the VideoNorms authors to request annotations and expects access soon. Annotation analysis remains pending receipt; no result here establishes its training value. The pooled CLIP diagnostic and selected cases further motivate representations that preserve interaction details and distinguish action choice from justification quality.
